\documentclass[12pt]{article}
\usepackage[T1]{fontenc}
\usepackage[utf8]{inputenc}
\usepackage{lmodern}
\usepackage{textcomp}
\usepackage{enumitem}
\usepackage{geometry}
\usepackage{graphicx}
\usepackage{amsmath, amssymb}
\geometry{margin=1in}
\begin{document}
\begin{center}
Engineering - Graduate Level Assessment 7 \
Total Marks: 40 \
Answer all questions.
\end{center}

\section*{Multiple Choice (10 marks)}
\begin{enumerate}[label=\arabic*.]
\item Convert the Fourier series f(t) = ^+\ensuremath{\infty}\ensuremath{\sum}\_k=-\ensuremath{\infty {(1) / (1 +jk)}}$e^jkt$ to trigonometric form.
\begin{enumerate}[label=(\alph*)]
    \item f(t) = 1 + 2 ^\ensuremath{\infty}\ensuremath{\sum}\_k=1 [\{coskt\} / \{1 + $k^2$\}]
    \item f(t) = ^\ensuremath{\infty}\ensuremath{\sum}\_k=1 [\{coskt+ k sinkt\} / \{1 + $k^2$\}]
    \item f(t) = 2 ^\ensuremath{\infty}\ensuremath{\sum}\_k=1 [\{coskt+ k sinkt\} / \{1 + $k^2$\}]
    \item f(t) = 1 + 2 ^\ensuremath{\infty}\ensuremath{\sum}\_k=1 [\{coskt+ k sinkt\} / \{1 + $k^2$\}]
    \item f(t) = 2 + 2 ^\ensuremath{\infty}\ensuremath{\sum}\_k=1 [\{coskt\} / \{1 + $k^2$\}]
\end{enumerate}
\item The inductance of linear time-varying inductor is given by L(t) = L\_0(t +tanht) and the current through it is given by i (t) =cos\ensuremath{\omegat} Find the voltage across this inductor.
\begin{enumerate}[label=(\alph*)]
    \item U(t) = L\_0 (1 - tanht)cos\ensuremath{\omegat} + \ensuremath{\omegaL}\_0 (t - $sech^2$ t)sin\ensuremath{\omegat}
    \item U(t) = L\_0 (1 + $sech^2$ t)cos\ensuremath{\omegat}+ \ensuremath{\omegaL}\_0 (t +tanht)sin\ensuremath{\omegat}
    \item U(t) = L\_0 (1 + $sech^2$ t)cos\ensuremath{\omegat}- \ensuremath{\omegaL}\_0 (t +tanht)sin\ensuremath{\omegat}
    \item U(t) = L\_0 (1 - $sech^2$ t)cos\ensuremath{\omegat}- \ensuremath{\omegaL}\_0 (t +tanht)sin\ensuremath{\omegat}
    \item U(t) = -L\_0 (1 + $sech^2$ t)sin\ensuremath{\omegat} + \ensuremath{\omegaL}\_0 (t + tanht)cos\ensuremath{\omegat}
\end{enumerate}
\item Compute the lowest natural frequency for an automobile enginespring. The spring is made of spring steel with 10 activecoils, mean diameter of 3.5 in. and wire diameter of 0.20 in.
\begin{enumerate}[label=(\alph*)]
    \item 15 cycles/sec
    \item 33 cycles/sec
    \item 22.5 cycles/sec
    \item 25.8 cycles/sec
    \item 18 cycles/sec
\end{enumerate}
\item Calculate the heat added when one mole of carbon dioxide is heated at constant volume from 540 to 3540 F.
\begin{enumerate}[label=(\alph*)]
    \item 28,500 Btu
    \item 34,650 Btu
    \item 40,000 Btu
    \item 25,000 Btu
    \item 45,000 Btu
\end{enumerate}
\item Both the ALU and control section of CPU employ which special purpose storage location?
\begin{enumerate}[label=(\alph*)]
    \item Multiplexers
    \item Memory Cache
    \item Registers
    \item Flip-Flops
    \item Buffers
\end{enumerate}
\end{enumerate}

\section*{True / False (10 marks)}
\textit{Answer True or False}
\begin{enumerate}[label=\arabic*.]
\setcounter{enumi}{5}
\item True or False: A DC machine commonly has both high shunt and armature resistance values to improve efficiency.
\item True or False: The Laplace transform of f(t) = sin(kt) is given by L\{sin(kt)\} = k / ($s^2$ - $k^2$) for t \textgreater{} 0.
\item True or False: The time-domain response v\_ab(t) for the given transform is -48 e^-3 t - 34 e^-2 t + 16.5 e^-4 t.
\item True or False: In the decimal number system, the most significant digit (MSD) is the last digit from left to right.
\item True or False: All oscillators listed, including the Hartley oscillator, produce sinusoidal waveforms under ideal conditions.
\end{enumerate}

\section*{Long Form Response (20 marks)}
\textit{Answer each question with a clear and complete explanation.}
\begin{enumerate}[label=\arabic*.]
\setcounter{enumi}{10}
\item Calculate and discuss the diffusion rate of water into dry air at 1 atm and 25^\ensuremath{\circC} from a test tube with a diameter of 12 mm and a length of 16 cm, including relevant equations and assumptions.
\item Discuss the process of determining the homogeneous solution for the voltage, v\_C(t), across a (4/3)F capacitor in a series RLC circuit with given parameters.
\end{enumerate}

\vfill
\noindent\textit{End of Paper}
\end{document}